% Please add the following required packages to your document preamble:
% \usepackage{multirow}
% \begin{table*}[!h]
% \centering 
% \setlength{\extrarowheight}{1pt}
% \begin{tabular}{l|lll|ll}
% \textbf{} & \textbf{Method} & \textbf{\begin{tabular}{l}
% Contrastive\\Combination\end{tabular}} & \textbf{Reduction} & \textbf{Pearson} & \textbf{Kendall-Tau} \\ \hline
% \multirow{6}{*}{\rotatebox{90}{eLife}} & ICL & - & - & -0.025 & -0.017 \\
%  % & \hspace{16pt}- prompt 2 & - & - &  &  \\ 
%  \cdashline{2-6}[1pt/3pt]
%  & \multirow{4}{*}{\begin{tabular}[c]{@{}l@{}}Control \\ Vectors\end{tabular}} & Diff & Avg & \textbf{0.302} & \textbf{0.205} \\
%  &  & Center & Avg & \textbf{0.302} & \textbf{0.205} \\
%  &  & Diff & PCA & 0.101 & 0.075 \\
%  &  & Center & PCA & 0.101 & 0.075 \\ \hline
%   %---------------------------------------------
% \multirow{6}{*}{\rotatebox{90}{PLOS}} & ICL & - & - & 2.7e-4 & -0.003 \\
%  % & \hspace{16pt}- prompt 2 & - & - &  &  \\ 
% \cdashline{2-6}[1pt/3pt]
%  & \multirow{4}{*}{\begin{tabular}[c]{@{}l@{}}Control \\ Vectors\end{tabular}} & Diff & Avg & 0.416 & 0.286 \\
%  &  & Center & Avg & \textbf{0.417} &  \textbf{0.287} \\
%  &  & Diff & PCA & 0.365 & 0.249 \\
%  &  & Center & PCA & 0.144 & 0.107 \\ \hline
%   %---------------------------------------------
% \multirow{6}{*}{\rotatebox{90}{Eureka}} & ICL & - & - & 0.055 & 0.0194 \\
%  % & \hspace{16pt}- prompt 2 & - & - &  &  \\ 
% \cdashline{2-6}[1pt/3pt]
%  & \multirow{4}{*}{\begin{tabular}[c]{@{}l@{}}Control \\ Vectors\end{tabular}} & Diff & Avg & 0.407 & 0.298 \\
%  &  & Center & Avg & \textbf{0.408} & \textbf{0.298} \\
%  &  & Diff & PCA & 0.062 & 0.043 \\
%  &  & Center & PCA & 0.168 & 0.1212 \\ \hline
%  %---------------------------------------------
% \multirow{6}{*}{\rotatebox{90}{SciNews}} & ICL & - & - &  0.084 & 0.025  \\
%  % & \hspace{16pt}- prompt 2 & - & - &  &  \\ 
% \cdashline{2-6}[1pt/3pt]
%  & \multirow{4}{*}{\begin{tabular}[c]{@{}l@{}}Control \\ Vectors\end{tabular}} & Diff & Avg & 0.116 & 0.075 \\
%  &  & Center & Avg & 0.268 &  0.216\\
%  &  & Diff & PCA & \textbf{0.395} & \textbf{0.315} \\
%  &  & Center & PCA & 0.081 & 0.074
% \end{tabular}
% \caption{Full set of results for all the Control Vector settings we tested. We report the Pearson and Kendall-Tau correlation of the specified level of readability with the Coleman-Liau readability index. We experiment with In-Context Learning (ICL) as our baseline. For the Control Vectors, we experiment with 2 Contrastive Combination methods (Diff and Center) and 2 Reduction methods (Avg. and PCA).}\label{tab:main-results}
% \end{table*}



% Please add the following required packages to your document preamble:
% \usepackage{multirow}
\begin{table*}[!ht]
\centering 
\footnotesize
\setlength{\extrarowheight}{1pt}
% Please add the following required packages to your document preamble:
% \usepackage{multirow}
\begin{tabular}{l|lll|lll}
\textbf{} & \textbf{Method} & \textbf{\begin{tabular}[c]{@{}l@{}}Contrastive\\ Combination\end{tabular}} & \textbf{Reduction} & \textbf{Pearson} & \textbf{Kendall Tau} & \textbf{BertScore F1} \\ \hline
\multirow{8}{*}{\rotatebox{90}{eLife}} & ICL & - & - & -0.025 & -0.017 & 0.847 \\
 & DSPy & - & - & 0.000 & 0.000 & 0.000 \\
 & QLoRa SFT & - & - & -0.115 & -0.164 & 0.847 \\
 & SFT & - & - & 0.033 & 0.035 & 0.780 \\  \cdashline{2-7}[1pt/3pt]
 & \multirow{4}{*}{\begin{tabular}[c]{@{}l@{}}Control \\ Vectors\end{tabular}} & Diff & Avg & 0.302 & 0.205 & 0.837 \\
 &  & Center & Avg & 0.302 & 0.205 & 0.803 \\
 &  & Diff & PCA & 0.101 & 0.075 & 0.836 \\
 &  & Center & PCA & 0.101 & 0.075 & 0.835 \\ \hline
\multirow{8}{*}{\rotatebox{90}{PLOS}} & ICL & - & - & 0.00024 & -0.003 & 0.848 \\
 & DSPy & - & - & 0.000 & 0.000 & 0.000 \\
 & LoRa SFT & - & - & -0.084 & -0.126 & 0.819 \\
 & SFT & - & - & 0.0074 & 0.004 & 0.719 \\  \cdashline{2-7}[1pt/3pt]
 & \multirow{4}{*}{\begin{tabular}[c]{@{}l@{}}Control \\ Vectors\end{tabular}} & Diff & Avg & 0.416 & 0.286 & 0.805 \\
 &  & Center & Avg & 0.417 & 0.287 & 0.805 \\
 &  & Diff & PCA & 0.365 & 0.249 & 0.838 \\
 &  & Center & PCA & 0.144 & 0.107 & 0.837 \\ \hline
\multirow{8}{*}{\rotatebox{90}{Eureka}} & ICL & - & - & 0.055 & 0.0194 & 0.853 \\
 & DSPy & - & - & 0.000 & 0.000 & 0.000 \\
 & LoRa SFT & - & - & -0.034 & -0.036 & 0.846 \\
 & SFT & - & - & 0.0078 & 0.0095 & 0.764 \\  \cdashline{2-7}[1pt/3pt]
 & \multirow{4}{*}{\begin{tabular}[c]{@{}l@{}}Control \\ Vectors\end{tabular}} & Diff & Avg & 0.407 & 0.298 & 0.840 \\
 &  & Center & Avg & 0.408 & 0.298 & 0.810 \\
 &  & Diff & PCA & 0.062 & 0.043 & 0.839 \\
 &  & Center & PCA & 0.168 & 0.1212 & 0.838 \\ \hline
\multirow{8}{*}{\rotatebox{90}{SciNews}} & ICL & - & - & 0.0302 & -0.004 & 0.840 \\
 & DSPy & - & - & 0.000 & 0.000 & 0.000 \\
 & LoRa SFT & - & - & -0.048 & -0.053 & 0.480 \\
 & SFT & - & - & -0.0062 & -0.0059 & 0.44 \\  \cdashline{2-7}[1pt/3pt]
 & \multirow{4}{*}{\begin{tabular}[c]{@{}l@{}}Control \\ Vectors\end{tabular}} & Diff & Avg & 0.075 & 0.000 & 0.827 \\
 &  & Center & Avg & 0.268 & 0.216 & 0.801 \\
 &  & Diff & PCA & 0.395 & 0.315 & 0.826 \\
 &  & Center & PCA & 0.081 & 0.074 & 0.826
\end{tabular}
\caption{Full set of results for all the Control Vector settings we tested for \texttt{meta-llama/Llama-3.1-8B-Instruct}. We report the Pearson and Kendall-Tau correlation of the specified level of readability with the Coleman-Liau readability index. We experiment with In-Context Learning (ICL), Supervised Finetuning (SFT), QLoRa/LoRa SFT, and DSPy Prompt Optimization (DSPy) as our baselines. For the Control Vectors, we experiment with 2 Contrastive Combination methods (Diff and Center) and 2 Reduction methods (Avg. and PCA).}
\label{tab:main-results-llama}
\end{table*}





% Please add the following required packages to your document preamble:
% \usepackage{multirow}
\begin{table*}[!ht]
\centering 
\footnotesize
\setlength{\extrarowheight}{1pt}
\begin{tabular}{l|lll|lll}
\textbf{} & \textbf{Method} & \textbf{\begin{tabular}[c]{@{}l@{}}Contrastive\\ Combination\end{tabular}} & \textbf{Reduction} & \textbf{Pearson} & \textbf{Kendall Tau} & \textbf{BertScore F1} \\ \hline
\multirow{7}{*}{\rotatebox{90}{eLife}} 
 & DSPy & - & - & \multicolumn{3}{c}{CONTEXT ERROR}  \\
 & LoRa SFT & - & - & \multicolumn{3}{c}{OOM} \\
 & SFT & - & - & \multicolumn{3}{c}{OOM} \\ 
 \cdashline{2-7}[1pt/3pt]
 & \multirow{4}{*}{\begin{tabular}[c]{@{}l@{}}Control \\ Vectors\end{tabular}} & Diff & Avg & 0.000 & 0.000 & 0.822 \\
 &  & Center & Avg & 0.132 & 0.206 & 0.808 \\
 &  & Diff & PCA & 0.315 & 0.223 & 0.821 \\
 &  & Center & PCA & 0.280 & 0.207 & 0.823 \\ \hline
\multirow{7}{*}{\rotatebox{90}{PLOS}} 
 & DSPy & - & - & \multicolumn{3}{c}{CONTEXT ERROR} \\
 & LoRa SFT & - & - & 0.000 & 0.000 & 0.000 \\
 & SFT & - & - & \multicolumn{3}{c}{OOM} \\ 
 \cdashline{2-7}[1pt/3pt]
 & \multirow{4}{*}{\begin{tabular}[c]{@{}l@{}}Control \\ Vectors\end{tabular}} & Diff & Avg & 0.000 & 0.000 & 0.832 \\
 &  & Center & Avg & 0.085 & 0.112 & 0.815 \\
 &  & Diff & PCA & 0.213 & 0.189 & 0.830 \\
 &  & Center & PCA & 0.100 & 0.120 & 0.831 \\ \hline
\multirow{7}{*}{\rotatebox{90}{Eureka}} 
 & DSPy & - & - & \multicolumn{3}{c}{CONTEXT ERROR} \\
 & LoRa SFT & - & - & -0.0017 & 0.0086 & 0.813 \\
 & SFT & - & - & 0.0031 & -0.011 & 0.798 \\ 
 \cdashline{2-7}[1pt/3pt]
 & \multirow{4}{*}{\begin{tabular}[c]{@{}l@{}}Control \\ Vectors\end{tabular}} & Diff & Avg & 0.000 & 0.000 & 0.906 \\
 &  & Center & Avg & 0.117 & 0.147 & 0.870 \\
 &  & Diff & PCA & 0.137 & 0.112 & 0.891 \\
 &  & Center & PCA & 0.121 & 0.098 & 0.899 \\ \hline
\multirow{7}{*}{\rotatebox{90}{SciNews}} 
 & DSPy & - & - & \multicolumn{3}{c}{CONTEXT ERROR} \\
 & LoRa SFT & - & - & -0.034 & -0.040 & 0.483 \\
 & SFT & - & - & \multicolumn{3}{c}{OOM} \\ 
 \cdashline{2-7}[1pt/3pt]
 & \multirow{4}{*}{\begin{tabular}[c]{@{}l@{}}Control \\ Vectors\end{tabular}} & Diff & Avg & 0.000 & 0.000 & 0.824 \\
 &  & Center & Avg & 0.141 & 0.177 & 0.809 \\
 &  & Diff & PCA & 0.116 & 0.110 & 0.822 \\
 &  & Center & PCA & 0.120 & 0.114 & 0.824
\end{tabular}
\caption{Full set of results for all the Control Vector settings we tested for \texttt{google/gemma-7b-it}. We report the Pearson and Kendall-Tau correlation of the specified level of readability with the Coleman-Liau readability index. We experiment with In-Context Learning (ICL), Supervised Finetuning (SFT), QLoRa/LoRa SFT, and DSPy Prompt Optimization (DSPy) as our baselines. For the Control Vectors, we experiment with 2 Contrastive Combination methods (Diff and Center) and 2 Reduction methods (Avg. and PCA). OOM means the baseline could not be run due to an out of memory error on our setup. UNSUPPORTED means that out of the box the baseline does not support \texttt{google/gemma-7b-it}.}
\label{tab:main-results-gemma}
\end{table*}



\begin{table*}[!ht]
\centering 
\footnotesize
\setlength{\extrarowheight}{1pt}
% Please add the following required packages to your document preamble:
% \usepackage{multirow}
\begin{tabular}{l|lll|lll}
\textbf{} & \textbf{Method} & \textbf{\begin{tabular}[c]{@{}l@{}}Contrastive\\ Combination\end{tabular}} & \textbf{Reduction} & \textbf{Pearson} & \textbf{Kendall Tau} & \textbf{BertScore F1} \\ \hline
\multirow{7}{*}{\rotatebox{90}{eLife}} 
 & DSPy & - & - & \multicolumn{3}{c}{CONTEXT ERROR} \\
 & LoRa SFT & - & - & 0.0011 & 0.0040 & 0.850 \\
 & SFT & - & - & 0.0012 & 0.000084 & 0.724 \\ 
 \cdashline{2-7}[1pt/3pt]
 & \multirow{4}{*}{\begin{tabular}[c]{@{}l@{}}Control \\ Vectors\end{tabular}} & Diff & Avg & 0.000 & 0.000 & 0.822 \\
 &  & Center & Avg & 0.132 & 0.206 & 0.808 \\
 &  & Diff & PCA & 0.315 & 0.223 & 0.821 \\
 &  & Center & PCA & 0.280 & 0.207 & 0.823 \\ \hline
\multirow{7}{*}{\rotatebox{90}{PLOS}} 
 & DSPy & - & - & \multicolumn{3}{c}{CONTEXT ERROR} \\
 & LoRa SFT & - & - & -0.00032& 0.0011 & 0.857 \\
 & SFT & - & - & 0.0027 & 0.0049 & 0.424 \\ 
 \cdashline{2-7}[1pt/3pt]
 & \multirow{4}{*}{\begin{tabular}[c]{@{}l@{}}Control \\ Vectors\end{tabular}} & Diff & Avg & 0.000 & 0.000 & 0.832 \\
 &  & Center & Avg & 0.085 & 0.112 & 0.815 \\
 &  & Diff & PCA & 0.213 & 0.189 & 0.830 \\
 &  & Center & PCA & 0.100 & 0.120 & 0.831 \\ \hline
\multirow{7}{*}{\rotatebox{90}{Eureka}} 
 & DSPy & - & - & \multicolumn{3}{c}{CONTEXT ERROR} \\
 & LoRa SFT & - & - & 0.027 & 0.035 & 0.766 \\
 & SFT & - & - & 0.000 & 0.000 & 0.696 \\ 
 \cdashline{2-7}[1pt/3pt]
 & \multirow{4}{*}{\begin{tabular}[c]{@{}l@{}}Control \\ Vectors\end{tabular}} & Diff & Avg & 0.000 & 0.000 & 0.906 \\
 &  & Center & Avg & 0.117 & 0.147 & 0.870 \\
 &  & Diff & PCA & 0.137 & 0.112 & 0.891 \\
 &  & Center & PCA & 0.121 & 0.098 & 0.899 \\ \hline
\multirow{7}{*}{\rotatebox{90}{SciNews}} 
 & DSPy & - & - & \multicolumn{3}{c}{CONTEXT ERROR} \\
 & LoRa SFT & - & - & 0.000 & 0.000 & 0.104 \\
 & SFT & - & - & 0.000 & 0.000 & 0.361 \\ 
 \cdashline{2-7}[1pt/3pt]
 & \multirow{4}{*}{\begin{tabular}[c]{@{}l@{}}Control \\ Vectors\end{tabular}} & Diff & Avg & 0.000 & 0.000 & 0.824 \\
 &  & Center & Avg & 0.141 & 0.177 & 0.809 \\
 &  & Diff & PCA & 0.116 & 0.110 & 0.822 \\
 &  & Center & PCA & 0.120 & 0.114 & 0.824
\end{tabular}
\caption{Full set of results for all the Control Vector settings we tested for \texttt{mistralai/Mistral-7B-Instruct-v0.3}. We report the Pearson and Kendall-Tau correlation of the specified level of readability with the Coleman-Liau readability index. We experiment with In-Context Learning (ICL), Supervised Finetuning (SFT), QLoRa/LoRa SFT, and DSPy Prompt Optimization (DSPy) as our baselines. For the Control Vectors, we experiment with 2 Contrastive Combination methods (Diff and Center) and 2 Reduction methods (Avg. and PCA). OOM means the baseline could not be run due to an out of memory error on our setup. LENGTH ERROR means the baseline could not be run due to the context window being too short.}
\label{tab:main-results-mistral}
\end{table*}



\section{Results}\label{sec:4-chapter-results}
We report the results of all the ablations in Table~\Cref{tab:main-results-llama}, ~\Cref{tab:main-results-gemma} and ~\Cref{tab:main-results-mistral}. Overall, we see that the baselines do not perform as well as the control vectors. We also see that some of the methods are either unsupported or cannot run due to out of memory errors. 

We see that Avg works the best as a reduction setting for all datasets except SciNews. This holds for both Diff and Center as the Contrastive Combination setting. For eLife, PLOS, and Eureka, Center achieves slightly better performance, although the difference is small. In general, the choice of reduction method makes the largest difference in terms of performance. While the Center-Avg settings perform best for 3 of the 4 datasets, SciNews achieving better performance with Diff-PCA indicates this setting is not universally optimal. Some tuning of the control vector settings may be necessary, although the overhead of this tuning is minimal compared to that of higher compute methods, such as finetuning, prompt optimization, and reinforcement-learning.

The failure of the training methods does not mean they did not learn how to summarize but rather they were unable to learn to control the readability of the summaries.  ICL, SFT, QLoRA/LoRA SFT, and prompt optimization perform poorly, with ICL and prompt optimization aligning with prior work showing they are insufficient for readability control \cite{ribeiro-etal-2023-generating}. Training-based methods also underperform, likely due to dataset differences (scientific texts vs. shorter news and fewer examples). Example outputs appear in Tables \ref{tab:sft-examples}–\ref{tab:dspy-examples}.


\begin{table*}[t]
\centering
\footnotesize
\begin{tabular}{l|p{.9\textwidth}}
\textbf{RL} & \textbf{SFT Output} \\\hline

1 & . These have to brain with for to. The, the brain, and the. This and to. These and, and the cells. The to a and to the heart that cells to is cells, are the their their. In, it the heart of the brain this. The. Next of are cells the brain the heart the and by. But the that the cells to to this the et al to the cells. These, and to the environment, a. We the mechanism. The, and in,. In the body and are is. This cells, the body, and of the heart to the body to of a a of, the brain. The, the. \\\hline

3 & brain a a of to the to, they. However. \\\hline

5 & , the that the cells to the of the the of of the visual in and of the brain are, the the body. The. In, but the time to the eyes the brain and the brain cells. The their are to the brain. Here, and the cells. \\\hline

7 & enerative their body to, the two the, in the body of, the their into. We from the cell of, the into the by their cells. \\\hline

9 & brain. In. \\\hline

\end{tabular}
\caption{Outputs generated by SFT model and the intended Readability Level (RL)}
\label{tab:sft-examples}
\end{table*}


\begin{table*}[t]
\centering
\begin{tabular}{l|p{.9\textwidth}}
\textbf{Readability Level} & \textbf{LoRA Output} \\\hline

1 & Small open reading frames (smORFs) are small stretches of DNA that code for proteins, but their translation into proteins has been difficult to detect. This is because they are too short to produce the high conservation scores that are used to predict which stretches of DNA code for proteins. To overcome this problem, Kondziella et al. have developed a new method called Poly-Ribo-Seq that uses a combination of biochemistry and sequencing to detect smORFs that are being translated into proteins. Using this method, they were able to detect smORFs in Drosophila cells that were not detected by other methods. These smORFs were found to be translated at similar levels to canonical proteins and had the same characteristics as canonical proteins. These findings suggest that many smORFs in Drosophila are being translated into proteins and may have important functions in the cell. Future experiments will be needed to confirm this and to determine the functions of these proteins. \\ \hline

3 & The discoidin domain receptors (DDRs) are a family of receptor tyrosine kinases (RTKs) that bind to collagen and play a key role in tissue development, cell migration, and disease progression. DDRs are composed of an extracellular region that binds to collagen, a transmembrane domain and a cytoplasmic region that contains the catalytic kinase domain. Unlike other RTKs, DDRs are constitutively dimerised in the absence of ligand and activate downstream signalling through a mechanism that is not well understood. Here, we show that DDR1 activation by collagen involves phosphorylation of the juxtamembrane and activation loop tyrosines in trans between dimers. Phosphorylation between DDR1 dimers requires specific contacts between the transmembrane domains but not between the extracellular domains. These findings define a unique activation mechanism for DDR1. \\ \hline

5 & Stem cell biology is a rapidly evolving field that is critical for our understanding of how tissues and organs form during development and how they are maintained in the adult. In the fruit fly Drosophila melanogaster, the lymph gland is a multi-lobed organ that produces blood cells throughout the larval stages. This organ contains a stem cell niche that maintains a pool of undifferentiated blood cell progenitors. While the adult hematopoietic hub is known to be the final site of hematopoiesis, the mechanism of blood cell formation in larvae has been largely unexplored. Here we identify a population of cells in the early larval lymph gland that express the Notch signaling pathway. These cells are multipotent, dividing and give rise to the domeless expressing blood cell progenitors that populate the lymph gland. Notch signaling is required for their maintenance and relies on Dpp signaling from the niche. These results reveal the presence of hematopoietic stem cells in the early larval lymph gland and establish Drosophila as a model for studying hematopoietic stem cell biology. The findings also shed light on the evolutionary conservation of hematopoietic stem cell biology between Drosophila and vertebrates. \\ \hline

7 & The malaria parasite Plasmodium causes a disease that kills more than a million people every year. To infect the human body, the parasite must first invade the liver cells. How this invasion takes place is still not well understood. Silvie et al. have now investigated how Plasmodium parasites enter human liver cells. The experiments show that two host proteins, CD81 and SR-BI, can be used by the parasite to invade the liver cells. The parasite uses different proteins to enter the liver cells of different species. The experiments also show that a parasite protein called P36 can bind to the host proteins and trigger the invasion of the liver cells. The experiments provide new insights into the invasion process and could ultimately help to develop new treatments for malaria. \\ \hline

9 & A large family of proteins called TRIM5 is thought to be a key part of the body's natural defense against viruses. These proteins prevent certain viruses from multiplying by binding to the proteins that form the outer shell of the virus. To find out how TRIM5 proteins work, researchers have made a detailed 3D image of a TRIM5 protein bound to a virus. The image shows that the protein forms a hexagonal pattern around the virus, with each hexagon made up of 12 smaller protein molecules. The next step will be to find out how the proteins interact with the virus to prevent it from multiplying. These findings could also help to develop new treatments for viruses. \\ \hline

\end{tabular}
\caption{Outputs generated by LoRA model}
\label{tab:lora-examples}
\end{table*}

\begin{table*}[t]
\centering
\begin{tabular}{l|p{.9\textwidth}}
\textbf{Readability Level} & \textbf{DSPy Output} \\\hline

1 & ""\\\hline

3 & ""\\\hline

5 & ""\\\hline

7 & ""\\\hline

9 & ""\\ \hline

\end{tabular}
\caption{Outputs generated by DSPy (empty because of length failure)}
\label{tab:dspy-examples}
\end{table*}
For most settings, using Center for the Contrastive Combination and Avg for the Reduction step performs best. Additionally, although the control vectors use eLife for the extraction dataset, the best control vectors generalize well to the other datasets, achieving similarly high correlations. This suggests that the resulting control vectors are in fact representing readability, rather than dataset specific information. 


\subsection{Baseline Discussion}\label{appendix:discussion}
The baselines were less simple to run than the control vectors leading to many four major issues: Resource Usage, Out of Memory Errors, and Lack of Support Out of the Box, Small Training Set. 
\paragraph{Resource Usage}
For DSPy, it also used way more resources than control vectors while performing worse. The DSPy method we use uses in context learning which forces the prompt to include extremely long examples. This could lead to some of the sources not being able to be summarized as the total length of the prompt was more than the max context length of the model. This was the main issue with DSPy for these smaller model. 

As for SFT and QLoRa/LoRa-SFT, it requires the more forward passes as finetuning is multiple epochs with the same amount of tokens. In addition, it requires expensive backward passes that require more memory. Even with this more computation, our method performs better while keeping training costs down. 
\paragraph{Out of Memory Errors}
Many of the training based methods required hefty amounts of resources and very long source documents led to out of memory errors with even QLoRa. While truncation could have prevented these errors, truncating at an arbitrary length could harm the overall plain language summary as this increases problems with reading comprehension, adds more hallucinations, and more \cite{ding2024fewertruncationsimprovelanguage}. While smarter engineering could have performed with other training optimizations, the amount of work to prevent these errors would have dwarfed the time it takes to set up the control vectors and could have led to even slower training times with the increase in communication time between devices and loading from RAM \cite{rajbhandari2020zeromemoryoptimizationstraining}. In addition, Gemma experienced more out of memory errors with less parameters  due to a different less memory efficient architecture. 
\paragraph{Lack of Support Out of the Box}
We could not run DSPy prompt optimization on the Gemma model due to an incompatibility between DSPy and the Gemma system prompt. It would have required modifying DSPy for specific set of models by adding additional adapters not available out of the box. 
\paragraph{Small Training Set}
We hypothesize that SFT performed worse in controllability than LoRA due to the number of trainable parameters relative to the size of the training dataset. Our SFT model was trained on only the 128 most readable examples, which likely provided insufficient signal to effectively update the large number of parameters involved in full fine-tuning. In contrast, LoRA introduces far fewer trainable parameters, making it better suited for small-data regimes. This is consistent with finding found in \cite{whitehouse-etal-2024-low} where they find that "in low-data scenarios, LoRA is a better alternative to full fine-tuning."

At first glance, this result may appear to contradict the findings of~\textcite{ribeiro-etal-2023-generating}. However, their experiments were conducted with a substantially larger training dataset. Additionally, their source documents consisted primarily of news articles, which are typically much shorter and less structurally complex than scientific articles, reducing the difficulty of the summarization task.



\subsection{Analysis of extraction dataset requirements.}\label{sec:other-datasets-exp}
We conduct experiments using the non-compliant datasets. We extract control vectors using the 128 most readable summaries from each train dataset, similar to the process in~\S\ref{sec:data-selection}. For evaluation, we randomly sample 248 examples from each test set. At inference, we use the control vectors extracted from each datasets' respective extraction set. This is an easier task than that presented above as it does not test generalization to other datasets. We use Center as our Contrastive Combination method and Avg as our reduction method, the best performing setting above. Results are below:
\begin{table}[h!]
\vspace{-3mm}
\centering\footnotesize
\begin{tabular}{l|rr}
Dataset & \multicolumn{1}{r}{Pearson} & \multicolumn{1}{c}{KT} \\ \hline
eLife & 0.302 & 0.205 \\
PLOS & 0.099 &  0.112 \\
Eureka & -0.195 & -0.137 \\
SciNews & 0.083 & 0.051
\end{tabular} 
\caption{Correlation of non-compliant datasets with specified readability level and CLI score.}
\label{tab:other-dataset-results}
\vspace{-4mm}
\end{table}

We find that the non-compliant datasets perform poorly. Eureka results in a moderate negative correlation, likely a result of the extraneous information present in the lay summaries, which can increase the readability scores. SciNews and PLOS have low correlation scores, indicating poor controlability. This contrasted with the high performance presented in~\Cref{tab:main-results-mini} provides evidence for the validity of our requirement framework.

\input{research-chapters/4-chapter/tables/latency-results}
\subsection{Latency testing}\label{appendix:latency-testing}

We conduct latency tests on the added overhead of using control vectors over standard LM inference. We randomly generate input IDs and measure the latency of the forward pass, as this is the portion of the generation pipeline where control vectors are used and have the potential to add overhead. We experiment with 3 batch sizes (1,4, 16) and 3 sequence lengths (512,1024,2048) for a total of 9 ablations. For each setting, we run 32 warm-up runs then measure the latency for 1024 inference calls. We report the mean latency in milliseconds per sample and per token. We additionally report the $\Delta$ latency and percent overhead. All the latency tests are run using Llama 3.1 8b Instruct (\texttt{meta-llama/Llama-3.1-8B-Instruct}) on a single NVIDIA A100 GPU, using half-precision floating-point. The results are reported in~\Cref{tbl:latency-results}. We find that using control vectors adds approximately a 5\% latency overhead when measured per sample or per token. We believe this difference is negligible when compared to other controllable generation methods, that require expensive training or hyper-parameter searches, or the additional token processing required for In-Context Learning. 


